\section{Discussion}
\label{sec:discussion}

\subsection{Implications for Diversity Measurement}

Layer Agreement separates two questions that a network-layer count can
otherwise combine. Network Top-2 share records concentration among observed
network transitions. Corridor Top-2 share records concentration among feasible
physical directions for the same candidate-bearing segments. In the observed
snapshot, the association between these quantities is stronger for island and
archipelagic service-facing units than for coastal-mainland units under all
three candidate-allocation rules.

This geographic difference is descriptive. One possible explanation is that
island paths reach submarine infrastructure through a smaller set of coastal
gateways, whereas coastal-mainland paths can include more terrestrial routing.
The current measurements do not test that mechanism, and exposure level alone
does not establish it.

The heterogeneous threshold classes also limit aggregate interpretations. A
broad network distribution projects to a concentrated corridor distribution
in 81 service-facing units, but 158 remain broad at both layers and some units
expand. Network-layer breadth therefore cannot be converted into a physical-
corridor breadth estimate using one universal contraction rate.

\subsection{Operational Scope}

Exposure and corridor concentration identify different intervention points,
but the measurements do not evaluate interventions. A closer service instance
or different interconnection could change how often paths enter the inter-
region candidate space. Access through distinct landing regions could change
the conditional corridor distribution. These are implications of the two
measured quantities, not effects estimated by this study.

Corridor concentration also does not measure cable multiplicity, landing-
station redundancy, correlated hazards, available rerouting capacity, or repair
processes. These dimensions are required before the reported feasible-corridor
distributions can be interpreted as operational resilience.

\subsection{Reproducibility and Ethics}

The study uses archived public RIPE Atlas measurements and infrastructure
metadata and reports aggregate country--measurement statistics. The pipeline
records measurement IDs, configuration values, mapping-resolution states,
filtering counts, and the inputs used to produce the tables and figures. The
paper's 80\% four-class rule must be distinguished from the repository's
tier-based auxiliary classification when artifacts are released.

The measurements do not target individual users and do not report individual
probe addresses. Candidate assignments are feasible hypotheses and should not
be interpreted as disclosure of an operator's actual cable use.
